% Options for packages loaded elsewhere
\PassOptionsToPackage{unicode}{hyperref}
\PassOptionsToPackage{hyphens}{url}
\PassOptionsToPackage{dvipsnames,svgnames,x11names}{xcolor}
\PassOptionsToPackage{space}{xeCJK}
\documentclass[
]{article}
\usepackage{xcolor}
\usepackage[margin=2.5cm]{geometry}
\usepackage{amsmath,amssymb}
\setcounter{secnumdepth}{-\maxdimen} % remove section numbering
\usepackage{iftex}
\ifPDFTeX
  \usepackage[T1]{fontenc}
  \usepackage[utf8]{inputenc}
  \usepackage{textcomp} % provide euro and other symbols
\else % if luatex or xetex
  \usepackage{unicode-math} % this also loads fontspec
  \defaultfontfeatures{Scale=MatchLowercase}
  \defaultfontfeatures[\rmfamily]{Ligatures=TeX,Scale=1}
\fi
\usepackage{lmodern}
\ifPDFTeX\else
  % xetex/luatex font selection
    \setmainfont[]{DejaVu Sans}
    \setmonofont[]{DejaVu Sans Mono}
  \ifXeTeX
    \usepackage{xeCJK}
    \setCJKmainfont[]{Noto Sans CJK SC}
          \fi
  \ifLuaTeX
    \usepackage[]{luatexja-fontspec}
    \setmainjfont[]{Noto Sans CJK SC}
  \fi
\fi
% Use upquote if available, for straight quotes in verbatim environments
\IfFileExists{upquote.sty}{\usepackage{upquote}}{}
\IfFileExists{microtype.sty}{% use microtype if available
  \usepackage[]{microtype}
  \UseMicrotypeSet[protrusion]{basicmath} % disable protrusion for tt fonts
}{}
\makeatletter
\@ifundefined{KOMAClassName}{% if non-KOMA class
  \IfFileExists{parskip.sty}{%
    \usepackage{parskip}
  }{% else
    \setlength{\parindent}{0pt}
    \setlength{\parskip}{6pt plus 2pt minus 1pt}}
}{% if KOMA class
  \KOMAoptions{parskip=half}}
\makeatother
\usepackage{longtable,booktabs,array}
\usepackage{calc} % for calculating minipage widths
% Correct order of tables after \paragraph or \subparagraph
\usepackage{etoolbox}
\makeatletter
\patchcmd\longtable{\par}{\if@noskipsec\mbox{}\fi\par}{}{}
\makeatother
% Allow footnotes in longtable head/foot
\IfFileExists{footnotehyper.sty}{\usepackage{footnotehyper}}{\usepackage{footnote}}
\makesavenoteenv{longtable}
\setlength{\emergencystretch}{3em} % prevent overfull lines
\providecommand{\tightlist}{%
  \setlength{\itemsep}{0pt}\setlength{\parskip}{0pt}}
\usepackage{bookmark}
\IfFileExists{xurl.sty}{\usepackage{xurl}}{} % add URL line breaks if available
\urlstyle{same}
\hypersetup{
  colorlinks=true,
  linkcolor={Maroon},
  filecolor={Maroon},
  citecolor={Blue},
  urlcolor={Blue},
  pdfcreator={LaTeX via pandoc}}

\author{}
\date{}

\begin{document}

\section{表示、逻辑、直觉的统一框架：形式指导构造 / 构造获得直觉 /
直觉搜索形式}\label{ux8868ux793aux903bux8f91ux76f4ux89c9ux7684ux7edfux4e00ux6846ux67b6ux5f62ux5f0fux6307ux5bfcux6784ux9020-ux6784ux9020ux83b7ux5f97ux76f4ux89c9-ux76f4ux89c9ux641cux7d22ux5f62ux5f0f}

\begin{quote}
\textbf{The Unified Framework of Representation, Logic, and Intuition}
\textbar{} 2026-08-11 纲领性文章 (Manifesto) --- 统摄三篇论文 P0/P1/P2,
定义统一架构与最终结论 姊妹篇:
P0《三通道等价》、P1《泛化作为归纳形式化》、P2《神经宏编译》
\end{quote}

\begin{center}\rule{0.5\linewidth}{0.5pt}\end{center}

\textbf{Repository:
\href{https://github.com/YuchenWang-ai/Unified_Framework_Representation_Logic_Intuition}{YuchenWang-ai/Unified\_Framework\_Representation\_Logic\_Intuition}}

\subsection{Abstract}\label{abstract}

We present a unified framework in which representation, logic, and
intuition are not three competing faculties but \textbf{three phases of
one generative cycle}. Formal syntax (representation) \emph{directs}
construction (logic): a well-formed, factorized grammar tells the system
what to build. Construction \emph{yields} intuition: definition-driven
unfolding, compiled into weights, becomes the reasoning fast path.
Intuition \emph{searches} form: the fast path explores combinations, and
what it finds is re-expressed in the unified formal language, feeding
construction again. The cycle is closed: \textbf{form guides
construction; construction earns intuition; intuition searches form.}
Three companion papers realize and verify each phase --- P1
(representation → generalization theory), P0 (three execution channels
equivalent), P2 (construction → intuition compilation) --- under one
architecture. The unified conclusion: once, under a unified formal
syntactic structure, a neural network clearly, intuitively, and
interpretably observes the difference and the characteristics of
intuition versus construction, we may with full confidence use the
unified formal language to direct logical construction, use precise
logical construction to earn intuition, and use agile intuitive
reasoning to exhaustively search form.

\begin{center}\rule{0.5\linewidth}{0.5pt}\end{center}

\subsection{1 统一框架 (The Unified
Framework)}\label{ux7edfux4e00ux6846ux67b6-the-unified-framework}

\textbf{三层一体.} 表示 (representation)、逻辑 (logic)、直觉 (intuition)
是同一生成循环的 三个阶段, 不是三种独立能力:

\begin{verbatim}
        形式指导构造                    构造获得直觉
  形式语法 ──────────► 逻辑构造 ──────────► 直觉推理
      ▲                                    │
      │                                    │
      └──────────── 直觉搜索形式 ────────────┘
               (直觉返回并再表达于统一形式)
\end{verbatim}

\textbf{三环闭路:}

\begin{enumerate}
\def\labelenumi{\arabic{enumi}.}
\tightlist
\item
  \textbf{形式指导构造 (Form directs construction)} ---
  统一的形式化语法结构 (表示层) 决定系统能构造什么; 语法正确 ⇒ 训练极快
  (定义精确少样本收敛);
\item
  \textbf{构造获得直觉 (Construction earns intuition)} ---
  定义驱动的展开 (构造通道) 经权重编译成为推理快路径 (直觉通道);
  构造正确 ⇒ 泛化极强 (OOD 零衰减);
\item
  \textbf{直觉搜索形式 (Intuition searches form)} ---
  快路径在新输入上跳步泛化, 搜索组合; 其结果再表达于统一形式, 反哺构造
  --- 泛化正确 ⇒ 直觉极强 (推理快路径, 生成/执行于统一形式).
\end{enumerate}

\textbf{统一架构的支撑:} token-native 六层投影表示 (B/C/S/G/P/A),
零硬编码求值器, 三执行通道 (构造/直觉/形式) 在统一形式下语义等价.

\subsection{2 三篇论文的定位
(Positioning)}\label{ux4e09ux7bc7ux8bbaux6587ux7684ux5b9aux4f4d-positioning}

\begin{longtable}[]{@{}
  >{\raggedright\arraybackslash}p{(\linewidth - 8\tabcolsep) * \real{0.2123}}
  >{\raggedright\arraybackslash}p{(\linewidth - 8\tabcolsep) * \real{0.1164}}
  >{\raggedright\arraybackslash}p{(\linewidth - 8\tabcolsep) * \real{0.3356}}
  >{\raggedright\arraybackslash}p{(\linewidth - 8\tabcolsep) * \real{0.1233}}
  >{\raggedright\arraybackslash}p{(\linewidth - 8\tabcolsep) * \real{0.2123}}@{}}
\toprule\noalign{}
\begin{minipage}[b]{\linewidth}\raggedright
论文
\end{minipage} & \begin{minipage}[b]{\linewidth}\raggedright
环节
\end{minipage} & \begin{minipage}[b]{\linewidth}\raggedright
核心命题
\end{minipage} & \begin{minipage}[b]{\linewidth}\raggedright
证据
\end{minipage} & \begin{minipage}[b]{\linewidth}\raggedright
LLM 技术映射
\end{minipage} \\
\midrule\noalign{}
\endhead
\bottomrule\noalign{}
\endlastfoot
\textbf{P1 泛化作为归纳形式化} & 形式→构造 & 定义完备 + 关系充分 ⇒
系统泛化; 因子化表示是原因 & EXP-C1/41/53/60/80 & 形式 = 注意力
(attention) \\
\textbf{P0 三通道等价} & 构造→直觉 (认知) & 构造=形式=直觉 语义等价;
直觉是编译的构造 & EXP-50/41/80 & 直觉 = flashing / 蒸馏 \\
\textbf{P2 神经宏编译} & 构造→直觉 (方法) &
把慢构造路径编译为快推理路径; 语义闭包不变 & EXP-50/51/80/90 & 构造 =
合成 CoT (synthetic CoT) \\
\end{longtable}

\textbf{技术映射 (成熟技术对应):} 形式 ↔ 注意力 (attention, token
间关系的定位); 构造 ↔ 合成思维链 (synthetic CoT, 精确的逐步逻辑构造);
直觉 ↔ 蒸馏的闪念 (flashing / distilled intuition, 快速推理上下文).

三篇共享同一表示、同一模型、同一权威判定口径, 从三个角度验证统一循环.

\subsection{3 相互引用
(Cross-References)}\label{ux76f8ux4e92ux5f15ux7528-cross-references}

\begin{itemize}
\tightlist
\item
  \textbf{P1 ← P0}: P1 §7 (语义路径编译) 引用 P0 的三通道等价 (R 命题);
  P1 的 G5 (直觉=推理快路径) 由 P0 认知地展开.
\item
  \textbf{P1 ← P2}: P1 §7 (编译) 由 P2 方法地实现 (S1-S5 五阶段); P1
  的''直觉难获得易跳步'' 由 P2 的编译机制解释.
\item
  \textbf{P0 ← P1}: P0 §2 (三通道共享定义源) 引用 P1 §4 (六层表示); P0
  的反柏拉图主义 由 P1 G0 的定义模糊证据支撑.
\item
  \textbf{P0 ← P2}: P0 §3 (直觉=编译的构造) 由 P2 的编译方法实现; P0 的
  C4 (获得与跳步解耦) 由 P2 的验证+回退机制支撑.
\item
  \textbf{P2 ← P0}: P2 §4.1 (零样本语义闭包) 引用 P0 的三通道一致; P2
  的验证 (S4) 使用 P0 的三通道交叉验证.
\item
  \textbf{P2 ← P1}: P2 §2 (定义链) 直接构建于 P1 §2; P2 的编译目标
  (直觉快路径) 是 P1 G5 的方法化.
\end{itemize}

\subsection{4 统一结论 (Unified
Conclusion)}\label{ux7edfux4e00ux7ed3ux8bba-unified-conclusion}

当我们在统一的形式化注意力语法结构下, 最终不可避免地、必然地通过神经网络
清晰、直观、可解释地观察到直觉与构造的边界时, 我们就可以有充足的信心:

\begin{itemize}
\tightlist
\item
  \textbf{用统一的注意力形式语言指导逻辑构造} (形式→构造: 注意力定位
  token 间关系);
\item
  \textbf{用精确的合成 CoT 逻辑构造获得 flashing 直觉} (构造→直觉:
  合成思维链 编译为快速推理);
\item
  \textbf{用蒸馏的直觉上下文穷举注意力形式} (直觉→形式:
  蒸馏快路径搜索组合, 再表达于统一形式).
\end{itemize}

形式、构造、直觉在统一的框架下相互适配,
组成再不分彼此的教学、训练、推理的 反馈迭代管线 --- 这也许是我们通往 ASI
的众多路径中, 相对较快的一条捷径.

\subsection{5 已验证的支柱 (Verified
Pillars)}\label{ux5df2ux9a8cux8bc1ux7684ux652fux67f1-verified-pillars}

\begin{longtable}[]{@{}
  >{\raggedright\arraybackslash}p{(\linewidth - 4\tabcolsep) * \real{0.2317}}
  >{\raggedright\arraybackslash}p{(\linewidth - 4\tabcolsep) * \real{0.6220}}
  >{\raggedright\arraybackslash}p{(\linewidth - 4\tabcolsep) * \real{0.1463}}@{}}
\toprule\noalign{}
\begin{minipage}[b]{\linewidth}\raggedright
支柱
\end{minipage} & \begin{minipage}[b]{\linewidth}\raggedright
结论
\end{minipage} & \begin{minipage}[b]{\linewidth}\raggedright
证据
\end{minipage} \\
\midrule\noalign{}
\endhead
\bottomrule\noalign{}
\endlastfoot
形式指导构造 (速度) & 定义精确 3 epochs 收敛; 定义模糊需 40 epochs &
EXP-53/61 \\
构造获得直觉 (泛化) & 因子化表示 ⇒ 外推零衰减 (2000位/进制60/匿名100\%)
& EXP-80/C1b \\
直觉搜索形式 (推理) & 快路径 7-9× 快; 82→5 token 跳步; 三通道等价 1.000
& EXP-50/51/90 \\
表示承载结构 & 符号置换不变 (0.987), 学关系非符号 & EXP-41 \\
\end{longtable}

\subsection{位置敏感性: 直觉与逻辑的分化 (2026-08-11
洞察)}\label{ux4f4dux7f6eux654fux611fux6027-ux76f4ux89c9ux4e0eux903bux8f91ux7684ux5206ux5316-2026-08-11-ux6d1eux5bdf}

\textbf{直觉位置不敏感, 逻辑位置敏感} --- 统一框架下两路径对 token
位置混乱的 鲁棒性不同:

\begin{longtable}[]{@{}
  >{\raggedright\arraybackslash}p{(\linewidth - 4\tabcolsep) * \real{0.1868}}
  >{\raggedright\arraybackslash}p{(\linewidth - 4\tabcolsep) * \real{0.2857}}
  >{\raggedright\arraybackslash}p{(\linewidth - 4\tabcolsep) * \real{0.5275}}@{}}
\toprule\noalign{}
\begin{minipage}[b]{\linewidth}\raggedright
路径
\end{minipage} & \begin{minipage}[b]{\linewidth}\raggedright
位置敏感性
\end{minipage} & \begin{minipage}[b]{\linewidth}\raggedright
证据
\end{minipage} \\
\midrule\noalign{}
\endhead
\bottomrule\noalign{}
\endlastfoot
直觉 (编译快路径) & 不敏感 (容忍排列混乱) & EXP-80: 100\% 匿名置换
1.000 \\
逻辑 (构造慢路径) & 敏感 (需确定位置-语义映射) & EXP-61h: 2 同义 token →
canonical 化; ≥3 → 崩 \\
\end{longtable}

\textbf{认知解释}: ``符号打乱不影响理解'' = 理解走直觉通道 (位置不敏感);
``集中逻辑时失去含义'' = 激活逻辑通道时位置歧义暴露。

\textbf{框架意义}: 这量化了''直觉 = 编译结构 (位置自由)'' vs ``逻辑 =
展开结构 (位置绑定)'' --- 三环闭路 (形式指导构造 / 构造获得直觉 /
直觉搜索形式) 中, 直觉搜索形式正是因为对位置不敏感,
才能在新组合上自由搜索。

唯一缺憾：这个结构是不对称的，意味着形式部分，可能存在某种可以精确解释的过程，在今日之我的能力范围之外。

\end{document}
